\documentclass[11pt]{article}

\usepackage[margin=1in]{geometry}
\usepackage{amsmath,amssymb}
\usepackage{booktabs}
\usepackage{enumitem}
\usepackage{hyperref}
\hypersetup{hidelinks}

\title{Scaling Book Exercises: Section 8}
\author{Lucas Yan}
\date{}

\begin{document}
\maketitle

\noindent
\textbf{Chapter:} Serving LLaMA 3-70B on TPUs\\
\textbf{Source scan:} \texttt{../handwritten/section\_08\_handwritten.pdf}\\
\textbf{Reference:} \url{https://jax-ml.github.io/scaling-book/applied-inference/}

\paragraph{TPU v5e constants.}
Each chip has 16 GB of HBM, bfloat16 peak compute
\[
C=1.97\times10^{14}\ \text{FLOPs/s},
\]
and HBM bandwidth
\[
W_{\mathrm{HBM}}=8.2\times10^{11}\ \text{bytes/s}.
\]

\section*{Question 1: LLaMA 3-405B lower bounds}

Let \(P=405\times10^9\) be the parameter count. A transformer forward pass
uses approximately two FLOPs per parameter per token, so
\[
\boxed{\text{FLOPs/token}\approx2P=810\times10^9}.
\]

On \(N\) TPU v5e chips, the ideal compute-bound latency is
\begin{align*}
T_{\mathrm{math}}
&=\frac{2P}{NC}\\
&=\frac{810\times10^9}
        {N(1.97\times10^{14})}\\
&=\boxed{\frac{4.11}{N}\ \text{ms/token}}.
\end{align*}

If parameter loading is the bottleneck, assume bfloat16 parameters. Their
total size is \(2P=810\) GB, so the ideal HBM-bound latency is
\begin{align*}
T_{\mathrm{parameters}}
&=\frac{2P\ \text{bytes}}{NW_{\mathrm{HBM}}}\\
&=\frac{810\times10^9}
        {N(8.2\times10^{11})}\\
&=\boxed{\frac{988}{N}\ \text{ms/token}}.
\end{align*}
Both are optimistic lower bounds: they assume perfect parallel scaling and
ignore collective latency and other runtime overheads.

\section*{Question 2: Serving LLaMA 3-8B at batch size 240}

Assume int8 weights, int8 KV caches, and an 8,192-token context. The relevant
LLaMA 3-8B dimensions are
\[
L=32,\qquad K=8,\qquad H=128,\qquad V=128{,}256.
\]

\subsection*{(a) Model parameters}

At one byte per int8 parameter,
\[
\boxed{M_{\mathrm{parameters}}=8\times10^9\ \text{bytes}=8\ \text{GB}}.
\]

\subsection*{(b) KV caches}

Each token stores one key and one value for every layer and KV head:
\[
M_{\mathrm{KV/token}}
=2LKH
=2(32)(8)(128)
=65{,}536\ \text{bytes}.
\]
Thus one 8,192-token sequence requires
\[
65{,}536(8{,}192)
=536{,}870{,}912\ \text{bytes}
\approx0.537\ \text{GB}.
\]
At batch size \(B=240\),
\[
\boxed{M_{\mathrm{KV}}\approx128.85\ \text{GB}}.
\]

\subsection*{(c) Peak working activations}

The largest unsharded working activation is approximately the logits tensor
\([B,V]\). Stored in bfloat16, it occupies
\[
M_{\mathrm{activations}}
\approx2BV
=2(240)(128{,}256)
=61{,}562{,}880\ \text{bytes},
\]
or
\[
\boxed{M_{\mathrm{activations}}\approx61.6\ \text{MB}}.
\]
This is small compared with the KV cache.

\subsection*{Smallest topology}

The total is approximately
\[
8+128.85+0.062=136.91\ \text{GB}.
\]
The raw capacity lower bound is therefore
\[
\left\lceil\frac{136.91}{16}\right\rceil=9\ \text{chips}.
\]
The next supported TPU v5e topology is a \(4\times4\) slice, so the smallest
usable topology is
\[
\boxed{4\times4=16\ \text{chips}}.
\]

\section*{Question 3: Serving LLaMA 3-405B}

Assume int8 weights, bfloat16 FLOPs, and a 128-Ki-token
\((S=131{,}072)\) context. LLaMA 3.1-405B has
\[
P=405\times10^9,\qquad L=126,\qquad K=8,\qquad H=128,
\qquad F=53{,}248.
\]

\subsection*{Memory}

The int8 model requires 405 GB. The KV cache per token is
\[
2LKH=2(126)(8)(128)=258{,}048\ \text{bytes}.
\]
One full-length sequence therefore requires
\[
258{,}048(131{,}072)
=33.82\times10^9\ \text{bytes}
=33.82\ \text{GB}.
\]
For batch size \(B\) and \(N\) chips, the memory constraint is
\[
405+33.82B\le16N.
\]
Even \(B=1\) requires at least
\[
\left\lceil\frac{405+33.82}{16}\right\rceil=28\ \text{chips},
\]
so the minimum supported topology is \(4\times8=32\) chips.

\subsection*{Latency model}

The ideal generation-step model is
\[
T_{\mathrm{step}}
=T_{\mathrm{KV}}+
\max\!\left(T_{\mathrm{parameters}},T_{\mathrm{math}}\right).
\]
In milliseconds,
\begin{align*}
T_{\mathrm{parameters}}
&=\frac{405\times10^9}{N(8.2\times10^{11})}(10^3)
=\frac{493.9}{N},\\
T_{\mathrm{KV}}
&=\frac{B(33.82\times10^9)}{N(8.2\times10^{11})}(10^3)
=\frac{41.25B}{N},\\
T_{\mathrm{math}}
&=\frac{2B(405\times10^9)}{N(1.97\times10^{14})}(10^3)
=\frac{4.11B}{N}.
\end{align*}
Compute overtakes parameter loading only near
\[
4.11B=493.9
\quad\Longrightarrow\quad
B\approx120.
\]
All feasible points below use \(B<120\), so
\[
T_{\mathrm{step}}\approx\frac{493.9+41.25B}{N}\ \text{ms}.
\]

The largest batch satisfying both memory and the 15 ms latency limit is
summarized below.

\begin{center}
\begin{tabular}{@{}rrrrrr@{}}
\toprule
\(N\) & Topology & \(B_{\mathrm{mem}}\) & \(B_{\mathrm{latency}}\)
& Chosen \(B\) & Step time \\
\midrule
32  & \(4\times8\)   & 3   & none & --- & \(>15\) ms \\
64  & \(8\times8\)   & 18  & 11   & 11  & 14.81 ms \\
128 & \(8\times16\)  & 48  & 34   & 34  & 14.82 ms \\
256 & \(16\times16\) & 109 & 81   & 81  & 14.98 ms \\
\bottomrule
\end{tabular}
\end{center}

Ignoring ICI communication, the corresponding aggregate throughputs are
approximately 743, 2,295, and 5,407 tokens/s. The naive roofline would
therefore choose 256 chips at batch size 81.

\subsection*{ICI roofline}

For ordinary decode tensor parallelism, the approximate useful sharding
limit is
\[
p_{\max}\approx\frac{F}{8B}.
\]
At \(N=256\) and \(B=81\),
\[
p_{\max}\approx\frac{53{,}248}{8(81)}\approx82<256,
\]
so that point is ICI-bound and its naive throughput estimate is not
achievable. At \(N=128\) and \(B=34\),
\[
p_{\max}\approx\frac{53{,}248}{8(34)}\approx196>128.
\]
Thus the simpler configuration that remains below this approximate ICI
roofline is
\[
\boxed{N=128\ (8\times16),\quad B=34,\quad
T_{\mathrm{step}}\approx14.82\ \text{ms}},
\]
with aggregate throughput
\[
\boxed{\frac{34}{0.01482}\approx2{,}295\ \text{tokens/s}}.
\]
A 256-chip deployment could instead run multiple smaller replicas; the
single-replica \(N=256,B=81\) calculation should not be treated as a valid
pure tensor-parallel operating point.

\subsection*{Theoretical minimum step time}

The most optimistic latency occurs at \(N=256,B=1\):
\begin{align*}
T_{\mathrm{parameters}}&\approx1.929\ \text{ms},\\
T_{\mathrm{KV}}&\approx0.161\ \text{ms},\\
T_{\mathrm{math}}&\approx0.016\ \text{ms}.
\end{align*}
Therefore
\[
T_{\min}
=0.161+\max(1.929,0.016)
=\boxed{2.09\ \text{ms/token}}.
\]
This is an ideal bandwidth lower bound; real ICI and fixed overheads make
the observed latency larger.

\end{document}
